\section{Discussão}
\label{sec:discussion}

Esta seção interpreta os resultados à luz das restrições operacionais e das escolhas metodológicas do estudo, enfatizando implicações práticas e limites de generalização.

Os resultados apresentados ao longo do artigo indicam que o preenchimento de lacunas em imagens de satélite deve ser tratado como um problema multicritério e dependente do contexto.
Essa afirmação pode parecer intuitiva, mas as evidências reunidas aqui permitem detalhá-la com maior precisão.
O comportamento dos métodos não depende apenas de sua formulação matemática ou de sua capacidade média de reduzir erro.
Ele depende também da estrutura espacial da cena, do nível de ruído no entorno observado e do custo computacional admissível para a aplicação final.

No bloco dos métodos clássicos, a principal implicação prática é que qualidade alta pode vir acompanhada de custo operacional elevado demais para determinadas aplicações.
Ordinary Kriging, \gls{RBF} e Thin Plate Spline aparecem repetidamente entre os métodos de melhor qualidade, mas não ocupam a mesma posição quando o tempo de inferência entra explicitamente na análise.
Esse resultado relativiza leituras puramente quantitativas baseadas em \gls{PSNR} ou \gls{SSIM}.
Em sistemas de produção, processamento em lote ou cenários de resposta rápida, a melhor técnica pode não ser aquela com maior qualidade média, e sim aquela que mantém desempenho competitivo sem comprometer a viabilidade computacional.

Esse ponto ajuda a explicar por que Exemplar-Based surge como método particularmente interessante no bloco clássico.
Ele não lidera todo o placar, mas se mantém próximo dos melhores em qualidade sem atingir o custo extremo de kriging ou spline.
Isso o posiciona como solução intermediária robusta em cenários nos quais a qualidade precisa ser alta, mas o custo não pode crescer sem controle.
Ao mesmo tempo, o artigo mostra que essa atratividade depende da cena.
Quando a lacuna ocorre em áreas cuja autossimilaridade é baixa, o ganho desse método tende a diminuir.

Outro resultado relevante é a ausência de uma relação simples e universal entre entropia local e dificuldade de reconstrução no bloco clássico.
As análises por correlação, degradação percentual e inclinação das curvas de \gls{PSNR} mostram que a entropia afeta os métodos de forma heterogênea.
Em alguns algoritmos, o aumento da entropia coincide com piora consistente de qualidade.
Em outros, o efeito é fraco, próximo de zero ou depende fortemente da métrica observada.
Isso significa que a entropia não deve ser interpretada como marcador universal de dificuldade com direção fixa.
Ela funciona melhor como variável moderadora, que reorganiza o regime de erro conforme o método e a escala espacial considerados.

Esse resultado é particularmente importante para o desenho de estudos futuros.
Se a entropia local fosse um marcador monotônico universal, bastaria estratificar o conjunto de dados em baixa, média e alta complexidade e esperar uma degradação regular em todos os métodos.
Os resultados mostram algo mais rico e mais desafiador.
Existe um regime intermediário de complexidade em que vários algoritmos se tornam especialmente instáveis.
Isso sugere que a dificuldade da reconstrução não cresce apenas com a diversidade tonal, mas também com a forma como essa diversidade se organiza espacialmente no entorno da lacuna.

No bloco de aprendizado profundo, a discussão assume outra forma.
A liderança da U-Net em qualidade absoluta é clara e consistente, o que reforça a efetividade da modelagem multiescala com conexões de atalho para esse problema.
No entanto, a análise de robustez mostra que qualidade absoluta e estabilidade relativa não são sinônimos.
\gls{GAN} e \gls{ViT} frequentemente degradam menos em termos percentuais quando o ruído aumenta.
Isso significa que a escolha da arquitetura depende do critério de desempenho considerado mais relevante para a aplicação final.

Essa distinção tem consequências práticas.
Se o objetivo é maximizar a qualidade final da reconstrução em um domínio semelhante ao do treinamento, a U-Net é a opção mais convincente entre as arquiteturas avaliadas.
Se o objetivo é obter comportamento mais previsível sob degradação do sinal observado, \gls{GAN} e \gls{ViT} passam a merecer atenção maior.
\gls{AE} e \gls{VAE}, embora inferiores em qualidade absoluta, continuam relevantes como referências de eficiência e simplicidade arquitetural.
Em outras palavras, o bloco de aprendizado profundo não produz apenas um vencedor.
Ele organiza um espaço de decisão com compromissos diferentes entre precisão, robustez e custo marginal.

Uma contribuição metodológica importante desta versão do manuscrito é a separação explícita entre os blocos clássico e profundo.
Essa escolha melhora a validade interpretativa do artigo porque reduz o risco de conclusões indevidas baseadas em comparação direta entre domínios experimentais heterogêneos.
Métodos clássicos foram analisados em regime multissensor com forte ênfase em custo e interpretabilidade.
Modelos de aprendizado profundo foram analisados no domínio Sentinel-2, com foco em diferenças arquiteturais internas.
Tratar esses dois contextos como se fossem intercambiáveis produziria uma narrativa mais simples, porém menos correta.

Essa decisão também afeta a forma como as evidências são apresentadas.
Ao separar tabelas de placar global, decomposição espectral, robustez e velocidade em blocos independentes, o artigo deixa mais claro o escopo de cada conclusão.
Essa reorganização não é apenas editorial.
Ela evita que a forma de apresentação produza uma inferência implícita de comparabilidade que os dados não sustentam com a mesma força.

Apesar da consistência interna dos resultados, o estudo possui limitações que precisam ser explicitadas.
A primeira é que o bloco clássico opera sobre um conjunto consolidado de avaliação balanceado, adotado para preservar comparabilidade e viabilidade computacional.
Isso fortalece a análise relativa entre métodos, mas restringe o alcance de generalização para todo o universo de recortes disponíveis no repositório.
A segunda limitação é a diferença de escopo entre os blocos clássico e profundo, que impede transformar a presença simultânea das duas famílias em argumento conclusivo de superioridade universal.

A terceira limitação diz respeito à modelagem estatística multivariada com múltiplas janelas de entropia.
Embora a regressão robusta seja útil para mostrar que a estrutura da cena permanece associada ao desempenho mesmo sob controle de método e ruído, as escalas de entropia são naturalmente correlacionadas entre si.
Isso enfraquece a interpretação isolada de coeficientes e recomenda cautela na leitura de efeitos marginais específicos.
Em outras palavras, o modelo é forte como evidência de associação multivariada, mas menos apropriado como instrumento de interpretação causal fina entre escalas.

Também é importante reconhecer que o artigo avalia qualidade de reconstrução, não impacto em tarefa final.
Uma técnica que preserve bem \gls{PSNR}, \gls{RMSE}, \gls{SAM} e \gls{ERGAS} ainda precisa ser validada em aplicações posteriores, como classificação temática, análise temporal ou estimativa biofísica.
Esse ponto delimita o alcance das conclusões sem diminuir sua relevância.
O estudo entrega um retrato consistente da reconstrução em si, que pode servir como base para avaliações orientadas a tarefa em trabalhos futuros.

Em síntese, a principal contribuição da discussão não está em declarar um único método vencedor.
Ela está em mostrar que a escolha metodológica correta depende da combinação entre estrutura local da cena, severidade da degradação, critério de desempenho e restrição operacional.
Ao explicitar essas dependências, o artigo oferece uma leitura mais útil e mais defensável do problema do que uma comparação direta e indiscriminada entre todas as técnicas avaliadas.
